\documentclass{leveldoc}
\DocUses{foundations,first_bend,twin_beams}
\DocFingerprint{c6e1d284c3444def}

\begin{document}

\section{What the level asks}\label{sec:asks}

The game's text:
\leveldescription

Two electrons ($q = \val{q}$, $m = \val{m}$) are launched from $(\val{xLaunch}, \val{yLaunch})$
along $+x$. One has kinetic energy $T = \val{Tslow}$ and must hit the one-cell spot
$\val{xS} \le x \le \val{xSlowHi}$, $\val{ySlowLo} \le y \le \val{ySlowHi}$. The other has
$T = \val{Tfast}$ and must hit $\val{xF} \le x \le \val{xFastHi}$,
$\val{yFastLo} \le y \le \val{yFastHi}$. Neither spot sets a direction or an energy window.
Each electron is tested only against its own spot. The player places up to three charges,
of either sign and of the magnitudes in the level file, inside the rectangle
$\val{xRlo} \le x \le \val{xRhi}$, $\val{yRlo} \le y \le \val{yRhi}$. A flight ends when
the electron touches a charge, leaves the world, enters its spot or runs out of time
(\xsecref{foundations}{sec:board}); $t_{\max} = \val{tmax}$.

The reference solution is the charge $Q = \val{Q}$ at $(\val{xC}, \val{yC})$, below the
rays. Both the electrons and the charge are negative, so $\kappa = qQ = \val{kappa}$ and
the force is repulsive: both rays bend upward. The slower electron bends more and meets
the higher spot. This document cites the orbit of one charge and derives the impulse that
makes the deflection fall as $1/T$.

\section{Setup}\label{sec:setup}

The level file, field by field (the appendix reproduces it):
\setuptable

In symbols: $c = \val{c}$, each electron has radius $a = 0$, and the charge has radius
$R = \val{R}$. The launches are
\[
  \gamma = 1 + \frac{T}{mc^2}, \qquad
  \abs{\vect{p}_0} = \sqrt{\frac{T^2}{c^2} + 2mT}.
\]
For the slower electron, $\gamma = \val{gammaSlow}$ and $\abs{\vect{p}_0} = \val{pSlow}$.
For the faster, $\gamma = \val{gammaFast}$ and $\abs{\vect{p}_0} = \val{pFast}$. Both
directions in the file are $(1, 0)$. The world is $-1 \le x \le \val{xWorld}$,
$-1 \le y \le \val{yWorld}$. The vertical offset of the charge below the rays is
$b = \val{b}$.

\section{Principles and model}\label{sec:model}

\begin{principles}
  \item The field of a charged sphere and the force on a non-polarizable particle,
    \xeqref{foundations}{eq:coulomb} and \xeqref{foundations}{eq:sphereforce}.
  \item The relativistic equation of motion \xeqref{foundations}{eq:eom}, the launch
    \xeqref{foundations}{eq:launch}, the energy \xeqref{foundations}{eq:energy} and the
    angular momentum \xeqref{foundations}{eq:angmom}. With one charge, the orbit
    \xeqref{first_bend}{eq:general}.
  \item The board's rules (\xsecref{foundations}{sec:board}): two shots are one puzzle,
    solved only when every shot arrives in its own detector.
  \item Newton's method and Simpson's rule (\xsecref{foundations}{sec:numerics}).
  \item The transverse impulse of a straight passage, \eqref{eq:impulse}. It estimates
    how the deflection depends on $T$. The arrival compared below is the orbit at the
    screen.
\end{principles}

With the charge at $\vect{r}_\mathrm{c}$ the force on either electron is
$\kappa(\vect{x}-\vect{r}_\mathrm{c})/\abs{\vect{x}-\vect{r}_\mathrm{c}}^3$, and the state
obeys \xeqref{foundations}{eq:eom} from its own launch until the first event. The two
flights do not interact. The motion stays in the plane (\xsecref{foundations}{sec:plane}).

\begin{assumption}{The charge is held in place}\label{as:fixed}
The charge does not recoil. That is a rule of the game (\xsecref{foundations}{sec:radiation}).
\end{assumption}

\begin{assumption}{No radiation}\label{as:radiation}
The radiated energy \xeqref{foundations}{eq:larmor} is not taken from the electron. Along
the two reference flights it is the fractions $\val{EradFracSlow}$ and $\val{EradFracFast}$
of the launch energy (§\ref{sec:hand}). The game makes the same approximation and reports
the same integral (\physref{8}).
\end{assumption}

\begin{assumption}{A point particle}\label{as:point}
$a = 0$, and the electron does not polarize, so \xeqref{foundations}{eq:sphereforce} is
exact. The closest approaches on the reference flights are $\val{rminSlow}$ and
$\val{rminFast}$, both far outside the sphere of radius $R$.
\end{assumption}

\begin{assumption}{Where the impulse applies}\label{as:impulse}
\eqref{eq:impulse} uses a straight path at constant speed, Newtonian $T = mv^2/2$, and a
flight that runs from infinity to infinity. Here each electron is launched at
$r_0 = \val{rSlow}$ from the charge, the screen cuts the orbit at $x = \val{xS}$, and
$\Gamma$ falls short of $1$ by $\val{gGap}$. The formula and the arrival angle differ by
$\val{impGapSlow}^\circ$ on the slow flight and by $\val{impGapFast}^\circ$ on the fast
one (§\ref{sec:impulse}). The comparison uses the orbit at the screen, which assumes none
of this.
\end{assumption}

\begin{figure}[htbp]
\centering
\begin{tikzpicture}
\begin{axis}[axis equal image, width=0.88\linewidth, xmin=-1.6, xmax=32.2, ymin=-1.6, ymax=22,
  xtick={0,5,...,30}, ytick={0,5,...,20}, tick label style={font=\small},
  xlabel={$x$}, ylabel={$y$}, label style={font=\small}]
  \draw[black!8] (0,0) grid (30,20);
  \draw[black!45, dashed] (-1,-1) rectangle (31,21);
  \draw[violet, dashed] (8,3) rectangle (17,17);
  \fill[green!50!black, fill opacity=0.35] (29,17) rectangle (30,18);
  \fill[green!50!black, fill opacity=0.35] (29,14) rectangle (30,15);
  \draw[green!40!black] (29,17) rectangle (30,18);
  \draw[green!40!black] (29,14) rectangle (30,15);
  \addplot[thin, red!70!black] table {pathhardslow.dat};
  \addplot[thin, red!70!black] table {pathhardfast.dat};
  \addplot[thin, black!40] table {pathweakslow.dat};
  \addplot[thin, black!40] table {pathweakfast.dat};
  \addplot[very thick, blue!70!black] table {pathslow.dat};
  \addplot[very thick, orange!80!black] table {pathfast.dat};
  \fill[red!70!black] (8,4) circle[radius=0.28];
  \node[font=\small, red!60!black, anchor=south] at (8,4.35) {$Q$};
  \draw[-{Stealth[length=6pt]}, thick] (0,10) -- (\rawval{arrowX},\rawval{arrowY});
  \fill (0,10) circle[radius=0.12];
  \node[font=\scriptsize, green!30!black, anchor=south] at (29.5,18.15) {$0.4$};
  \node[font=\scriptsize, green!30!black, anchor=north] at (29.5,13.85) {$0.6$};
\end{axis}
\end{tikzpicture}
\caption{The slower electron (blue) meets the upper spot and the faster electron (orange)
meets the lower spot. The dashed rectangle is where a charge may be placed. A weaker
charge at the same point (grey; §\ref{sec:others}) leaves both rays short of their spots.
A stronger charge (red) drives the slower electron into the ceiling and the faster
electron out above its own spot.}
\label{fig:board}
\end{figure}

\section{Derivation}\label{sec:derivation}

\subsection{The orbit of one electron}\label{sec:orbit}

Take polar coordinates about the charge and let $u = 1/r$. With
$\sigma = \operatorname{sign} L_z$, $L = \abs{L_z}$, $\theta = \sigma(\varphi - \varphi_0)$
and $u'_0 = -p_{r,0}/L$,
\begin{equation}\label{eq:orbit}
  u(\theta) = B + (u_0 - B)\cos\Gamma\theta + \frac{u'_0}{\Gamma}\sin\Gamma\theta,
\end{equation}
where $\Gamma^2 = 1 - \kappa^2/(c^2 L^2)$ and $B = -\kappa W/(c^2 L^2 \Gamma^2)$. This is
\xeqref{first_bend}{eq:general}. The time element and the momentum are
\xeqref{first_bend}{eq:dtdth} and \xeqref{first_bend}{eq:momentum}.

Both launches point along $+x$ from the same point, so the radius $r_0$ and the direction
$\varphi_0$ are shared, and $u'_0$ is shared as well: $p_{r,0}$ and $L$ both scale with
$\abs{\vect{p}_0}$. The radial momentum points toward the charge, $u'_0 = \val{duSlow}$,
so the launch is not an apsis. $L_z < 0$ on both flights. Repulsion from a charge below
the ray pushes each electron up, and the slower one, with the smaller $L$, is turned
through a larger angle.

\subsection{The screen}\label{sec:face}

The position on the orbit is
\begin{equation}\label{eq:position}
  x(\theta) = x_\mathrm{c} + \frac{\cos\varphi}{u(\theta)}, \qquad
  y(\theta) = y_\mathrm{c} + \frac{\sin\varphi}{u(\theta)}, \qquad
  \varphi = \varphi_0 + \sigma\theta,
\end{equation}
and $\dd x/\dd\theta$ is the derivative used in \xsecref{twin_beams}{sec:face}. Which
boundary comes first is read off the orbit. Walk $\theta$ from $0$ in steps of $10^{-3}$.
On each chord, in the board's order (\xsecref{foundations}{sec:board}): a radius
$r \le R$ is a collision; a point outside the world ends the flight through the side the
chord reaches first; a point inside that electron's own spot ends it through the side the
chord reaches first. The root is polished by Newton's method
\xeqref{foundations}{eq:newton}, started at the middle of the crossing chord. The face
coordinate of the root is set to the face value.

On both reference flights the first event is the screen face $x = x_S = \val{xS}$. The
root of
\begin{equation}\label{eq:face}
  f(\theta) = x(\theta) - x_S = 0
\end{equation}
is the arrival, and the ordinate there decides whether it lies in the spot.

\subsection{The time and the radiated energy}\label{sec:time}

The time along the orbit is the integral of \xeqref{first_bend}{eq:dtdth}. Level~1
evaluates it in closed form when $u'_0 = 0$. Here $u'_0 \ne 0$, and the number compared
below is Simpson's rule \xeqref{foundations}{eq:simpson} with $\val{nSimp}$ panels, as in
\xsecref{twin_beams}{sec:time}. Halving the panel count estimates the error as
$(S_n - S_{2n})/15$.

The radiated energy is the same rule on the power. The acceleration is
\xeqref{first_bend}{eq:acceleration} and the power is \xeqref{foundations}{eq:larmor}.
The integral uses $\val{nRad}$ panels; the uncertainty is the relative step-doubling error.

\subsection{The impulse, and the factor $1/T$}\label{sec:impulse}

Suppose instead that the electron held the straight line $y = y_0$ at a constant speed
$v$, with the charge a distance $b$ below that line, and that the line ran from
$x = -\infty$ to $x = \infty$. The transverse force is $\kappa b/r^3$ with
$r^2 = (x - x_\mathrm{c})^2 + b^2$, and $\dd t = \dd x/v$, so the kick is
\[
  \Delta p_y = \frac{\kappa b}{v}\int_{-\infty}^{\infty}\frac{\dd x}{(x^2 + b^2)^{3/2}}.
\]
The antiderivative is $x/(b^2\sqrt{x^2 + b^2})$: differentiating it brings down
$b^2/(x^2 + b^2)^{3/2}$ over $b^2$. The definite integral equals $2/b^2$, and
\[
  \Delta p_y = \frac{2\kappa}{vb}.
\]
In Newtonian mechanics $T = mv^2/2$ and $p = mv$, so $vp = 2T$. For a small kick the
ray turns through
\begin{equation}\label{eq:impulse}
  \theta \approx \frac{\Delta p_y}{p} = \frac{\kappa}{Tb}.
\end{equation}
The angle is proportional to $1/T$. Two energies on the same charge and the same offset
therefore turn in the ratio of the energies, and a screen far downstream separates them
by that ratio.

Read with $b = \val{b}$ and with $\kappa = \val{kappa}$, \eqref{eq:impulse} gives
$\val{impSlow}^\circ$ for the slow electron and $\val{impFast}^\circ$ for the fast one.
The momenta at the screen make $\val{angSlow}^\circ$ and $\val{angFast}^\circ$ with $+x$.
\ref{as:impulse} records the difference. The rises of the two orbits from the launch line,
$y - y_0$, are in the ratio $\val{ratioDy}$, against $T_\mathrm{fast}/T_\mathrm{slow}
= \val{ratioT}$. The formula is why the higher spot belongs to the lower energy. The
arrival is \eqref{eq:position} at the root of \eqref{eq:face}.

\section{The reference solution by hand}\label{sec:hand}

\begin{procedure}
  \item\label{st:launch} \emph{The launches} (\xeqref{foundations}{eq:launch}). Slow:
    $\gamma = \val{gammaSlow}$, $\abs{\vect{p}_0} = \val{pSlow}$. Fast:
    $\gamma = \val{gammaFast}$, $\abs{\vect{p}_0} = \val{pFast}$. Both aim along $+x$.
  \item\label{st:orbit} \emph{The slow orbit} \eqref{eq:orbit}, about $(\val{xC}, \val{yC})$.
    The launch radius is $r_0 = \val{rSlow}$ at $\varphi_0 = \val{phiSlow}$, and
    \begin{align*}
      L_z &= \val{LzSlow}, & W &= \val{WSlow}, & \Gamma &= \val{GammaSlow}, \\
      B &= \val{BSlow}, & u_0 &= \val{uSlow}, & u'_0 &= \val{duSlow}.
    \end{align*}
    The fast orbit has the same $r_0$, $\varphi_0$ and $u'_0$, with
    $L_z = \val{LzFast}$, $W = \val{WFast}$, $\Gamma = \val{GammaFast}$ and
    $B = \val{BFast}$.
  \item\label{st:face} \emph{The first event}, by the walk of §\ref{sec:face}. Both
    electrons meet $x = \val{xS}$ before a collision or a world edge. Points of the two
    orbits:
    \begin{center}\small\val{samples}\end{center}
    Newton's method on \eqref{eq:face} for the slow electron, started at
    $\theta = \val{thGuess}$:
    \begin{center}\small\val{newton}\end{center}
    After $\val{newtonSteps}$ corrections the residual is $\val{newtonResSlow}$, and
    $\theta = \val{thetaSlow}$. The ordinate is $\val{ySlow}$, inside the upper spot.
    The same method on the fast electron gives $\theta = \val{thetaFast}$ with residual
    $\val{newtonResFast}$, and the ordinate $\val{yFast}$ lies inside the lower spot.
    The closest approaches are at $\theta = \val{thMinSlow}$ and $\theta = \val{thMinFast}$,
    where $r = \val{rminSlow}$ and $r = \val{rminFast}$.
  \item\label{st:time} \emph{The time and the momentum.} Simpson's rule with
    $\val{nSimp}$ panels gives $t = \val{tSlow}$ and $t = \val{tFast}$, uncertain by
    $\val{terrSlow}$ and $\val{terrFast}$, both far below $t_{\max}$. The momenta are
    $(\val{pxSlow}, \val{pySlow})$ and $(\val{pxFast}, \val{pyFast})$, at
    $\val{angSlow}^\circ$ and $\val{angFast}^\circ$ to $+x$.
  \item\label{st:impulse} \emph{The impulse, as a check of the layout.}
    \eqref{eq:impulse} gives $\val{impSlow}^\circ$ and $\val{impFast}^\circ$. The rises
    from $y = \val{yLaunch}$ are $\val{dySlow}$ and $\val{dyFast}$, in the ratio
    $\val{ratioDy}$.
  \item\label{st:radiation} \emph{The radiated energy}, with $\val{nRad}$ panels:
    $\val{EradSlow}$ and $\val{EradFast}$, uncertain by the relative amounts
    $\val{radErrSlow}$ and $\val{radErrFast}$.
\end{procedure}

\begin{result}[The reference flights]
The slower electron enters its spot through $x = \val{xS}$ at $(\val{xSlow}, \val{ySlow})$,
at $t = \val{tSlow}$. The faster electron enters its spot through the same face at
$(\val{xFast}, \val{yFast})$, at $t = \val{tFast}$. Both ordinates lie in their spots.
The level is solved.
\end{result}

The orbits are the closed form \eqref{eq:orbit}. One Newton table is printed; the fast
root is the same method. The quadratures are specified in §\ref{sec:time}.

\section{Comparison with the game}\label{sec:comparison}

\comparisontable

The orbits are exact. The time quadratures are uncertain by $\val{terrSlow}$ and
$\val{terrFast}$ on the reference, and on the other flights by at most
$\val{HardSlowErr}$, all far below $10^{-12}$. The tolerance on a state is $10^{-9}$,
ten times the accuracy of a verified flight (\xsecref{foundations}{sec:verification}).
The tolerance on a radiated energy is $10^{-4}$ relative. Both were set from those
estimates. The radiated integrals are uncertain by at most $\val{WeakFastRadErr}$
relative.

The faster electron and the two weaker-charge flights stop on a vertical edge, and the
slower electron of the stronger charge stops on the ceiling, so those face coordinates
agree exactly. The largest absolute difference in the table is $\val{cmpHardSlowTDiff}$,
the time at which that stronger charge drives the slow electron into the ceiling. It sits
at the game's accuracy, above the quadrature error. The largest relative difference
among the radiated energies is $\val{cmpWeakFastRadDiff}$.

\section{Other solutions}\label{sec:others}

Two further charges at the same node $(\val{xC}, \val{yC})$, each flown by the same orbit
and compared above. Both miss the spots.

\begin{enumerate}
  \item $Q = \val{WeakQ}$, so $\kappa = \val{WeakK}$. Both electrons leave the world
    through $x = \val{xWorld}$. The slow one is then at $y = \val{WeakSlowY}$, below its
    spot, at $t = \val{WeakSlowT}$ and at $\val{WeakSlowAng}^\circ$. The fast one is at
    $y = \val{WeakFastY}$, below its spot, at $\val{WeakFastAng}^\circ$. The impulse
    \eqref{eq:impulse} has three quarters of the reference kick, and neither rise reaches
    its spot.
  \item $Q = \val{HardQ}$, so $\kappa = \val{HardK}$. The slow electron meets the ceiling
    $y = \val{yWorld}$ at $x = \val{HardSlowX}$, at $t = \val{HardSlowT}$ and at
    $\val{HardSlowAng}^\circ$, above its spot. The fast electron leaves through
    $x = \val{xWorld}$ at $y = \val{HardFastY}$, above its spot, at $\val{HardFastAng}^\circ$.
    One and a half times the reference kick carries the slow ray into the ceiling before
    the screen can catch it.
\end{enumerate}

The reference works because $\kappa = \val{kappa}$ and $b = \val{b}$ put the two rises,
in the ratio of the energies, one in each one-cell spot. A weaker charge at that point
undershoots both spots. A stronger charge overshoots them.

\section{What the game computes}\label{sec:engine}

The game integrates \xeqref{foundations}{eq:eom} for each electron separately, with its
eighth-order Runge--Kutta method at the tolerances $10^{-10}$ and $10^{-12}$ (\physref{5}).
It finds each crossing on its interpolant (\physref{6}) and verifies the outcome by
comparing its two runs (\physref{7}). The model is this document's, including
\ref{as:radiation} (\physref{8}). The impulse \eqref{eq:impulse} is not part of what the
game computes. The agreement in §\ref{sec:comparison} is between the orbit
\eqref{eq:orbit} and that integration.

\section{Further reading}\label{sec:reading}

\begin{itemize}
  \item \levelref{first_bend}: the orbit \eqref{eq:orbit}, the time element and the
    momentum. There the launch is an apsis, so the time is a closed form.
  \item \levelref{twin_beams}: the same orbit with $u'_0 \ne 0$, and Simpson's rule for
    the time. There the two rays are mirrors. Here they differ by their energy, and the
    separation is \eqref{eq:impulse}.
  \item \xsecref{foundations}{sec:board}: several shots form one puzzle. Each electron
    has its own spot.
\end{itemize}

\levelfileappendix

\end{document}
